\subsection{Ternary term operations}
Some families of negative source bits specify a single positive profile
bit. Given $A\subseteq V$ and an owner $r$, write $e(A,r)$ for the
coloring assigning $r$ on $A$ and leaving all other inputs uncolored.
For every core operation $H$,
\begin{equation}\label{core:eq:total-monochromatic-source}
 \bigl(\forall j\notin A,\ P_j^H(e(A,r))=0\bigr)
 \quad\Longleftrightarrow\quad \Phi_r(H)(A)=1 .
\end{equation}
Indeed, a positive profile forces the $r$ singleton at every one of
the displayed source probes. Conversely give the inputs outside $A$
the ambiguous menu $I\setminus\{r\}$. If the output contains a foreign
winner, source decomposition traces it to some $j\notin A$ with
$P_j^H(e(A,r))=1$. In particular a zero on a total monochromatic
passive coloring is a singleton profile demand by itself.

When the number of owners is part of the input, even these monochromatic
samples give a coNP-complete feasibility problem. The input lists demands
$r:A$ by the corresponding group of negative queries in
\eqref{core:eq:total-monochromatic-source}, and requests one fixed
source at the uncolored positive query. By \cref{core:thm:profiles},
the demands can coexist exactly when each selection of at most one
coalition per owner has a common input. The least profiles generated
by the demands and $V$ then have a greatest completion; it supplies
the marked source since no input of its marked probe is a singleton.
A selection with empty intersection certifies failure.

For hardness take the variables of a nonempty 3CNF as owners, padding
to three if needed, and its clauses as inputs. For each variable and
each of its two truth values require the coalition of clauses that this
choice alone does not satisfy. An assignment satisfies the formula
exactly when the corresponding coalitions have empty intersection.
Any partial selection with empty intersection extends to such an
assignment. The negative queries have polynomial total size, giving a
reduction from unsatisfiability. Here the number of inputs varies;
the bounded-arity test below applies even when the owner set varies.

\begin{corollary}[Joint guard recognition on three inputs]
\label{core:cor:ternary-guard-recognition}
For any finite $K\geq3$, consider a nonempty finite input set and
observer-agreeing antitone source predicates with at least one positive
bit. For each $r$, form the family
\[
 \mathcal W_r=\{A\subseteq V:\ P_j(e(A,r))=0\text{ for all }j\notin A\}.
\]
Prescribe $\Phi_r(A)=1$ for every inclusion-minimal
$A\in\mathcal W_r$. Discard any negative probe $(j,t)$ with
$t_r\in\mathcal W_r$ for some $r$. Suppose each remaining minimal
zero uses each color at most once. The joint predicates are realizable precisely
when the profile and marked reachability tests of
\cref{core:thm:repeated-blocker-budgets} hold for each maximal
positive source bit. In particular this is an intrinsic test for
\emph{all} core operations on at most three inputs, for any number of
owners $K\geq3$. On explicitly given guard tables it takes polynomial
time at each fixed arity in this class, even with $K$ part of the input.
At arity three it suffices to test rainbow selections of at most four
demands.
\end{corollary}
\begin{proof}
The families $\mathcal W_r$ are upward closed by antitonicity.
The tests give, for each maximal positive bit, a core word satisfying
that bit, the injectively colored zeros, and the prescribed singleton
profiles. A discarded probe returns $\{r\}$ because its $r$-coalition
is $t_r$; its marked input contributes no foreign winner.
Thus the word satisfies every minimal zero. Union over the positive witnesses and use
antitonicity. A realizing term conversely witnesses all the tests by
\eqref{core:eq:total-monochromatic-source}. With at most two passive
inputs, a repeated color fills the entire passive set: its complement
contains only the marked input. Thus these cases exhaust all minimal zeros.

For the complexity assertion put $n=|V|$. A rainbow selection with
empty intersection already fails the common-source condition, and at
most $n$ of its coalitions witness this. Otherwise, at each vertex
outside the common intersection keep up to two distinct demands whose
coalitions exclude it. The retained selection has the same intersection
and the same seed sets $L_r$: membership of a vertex in $L_r$ says
that some demand of an owner other than $r$ excludes it. These data
give the same projection and reachability tests on every subfamily
of the negative rows. Thus at most $2(n-1)$ demands suffice in the
nonempty-intersection case. There are at most $2^nK$ demands to choose
from; each individual reachability test takes polynomial time by
\cref{core:thm:repeated-blocker-budgets}.
\end{proof}

\begin{corollary}[Uniform ternary interpolation]
\label{core:cor:ternary-uniform-interpolation}
For every $K\ge3$, a ternary operation on winner menus is core if it
agrees with some core term on every set of at most ten valuations.
More generally, let $n\ge3$ and consider an operation reconstructed
from source guards satisfying the hypotheses of
\cref{core:cor:ternary-guard-recognition}. If it is not core, at most
$1+n\min(K,n)$ valuations witness this.
\end{corollary}
\begin{proof}
Put $b=\min(K,n)$. A prescribed demand $r:A$ costs one valuation:
give $A$ the $r$ singleton and all other inputs $I\setminus\{r\}$.
The reconstructed operation returns $\{r\}$ by the definition of
$\mathcal W_r$. Agreement there forces that demand on an interpolant.

Fix a positive bit whose bounded realization fails. In a failed
marked test, \cref{core:cor:input-guard-cuts} gives at most $nb$
negatives, plus the positive probe. In a failed rainbow test, an
empty intersection already has at most $n$ demand witnesses.
Otherwise let $C$ be the common intersection after compression in
\cref{core:cor:ternary-guard-recognition}. At a closed failing
subfamily, keep at most $b$ reachability witnesses per head in $C$
as in \cref{core:cor:input-guard-cuts}. Outside $C$ no projection
is eligible, and each vertex belongs to all $L_r$ except possibly
one, since some demand excludes it. Keep just the witness first
reaching it at that owner if needed. This preserves every reachability
set, every tail remains reached, and every eligible projection still
has a negative head. With the at most $2(n-|C|)$ demand valuations,
the total is at most $b|C|+3(n-|C|)\le bn$.

For a failed gap test arising from a demand $r:A$ and the positive
mark, every input outside $A$ is in all foreign seed sets.
The same argument uses at most $b|A|+n-|A|$ negatives. The demand
valuation and the positive probe enforce both signs of the gap.
If $A\subsetneq V$ the total is at most $bn$ since $b\ge3$;
the demand at $V$ needs no valuation and its gap follows from the mark.

Finally an arbitrary ternary operation locally interpolable on four
rows already has the required source representation. Read its bits
on the canonical probes of \cref{core:thm:source-guards}. Interpolation
on two probes gives antitonicity and agreement across observer overlaps.
For any valuation and winner, take that valuation together with the
at most three source probes of the inputs containing the winner. A
common term gives source decomposition on those rows, so the proposed
bits reconstruct the operation. Nonemptiness of its outputs then
ensures a positive bit. Every ternary candidate now falls under
\cref{core:cor:ternary-guard-recognition}, and $1+3b=10$.
\end{proof}

There is also a bound on the number of owner labels needed to test a
ternary operation given by such guards. It is core exactly when every
restriction to a palette $J\subseteq I$ with $3\le |J|\le20$ is core.
A term restricts by replacing each outside owner's choice by union.
Conversely, in the failure certificate above each of the at most ten
valuations asks for a source bit using at most two passive colors, or
for a singleton profile at one owner. Put these labels in $J$ and pad
to three if necessary. Each source query can be read using any mark
in $J$ missing from its coloring. The restricted operation gives the
same bits by guard reconstruction; its singleton demands are read by
the same monochromatic test with $J$ in place of $I$. A realizing word
over $J$, read over $I$, would therefore meet all the original
constraints. This contradicts the certificate.

\paragraph{Unanimous profiles at any number of owners.}
On three inputs, unanimity forces every source bit with at most one
passive color to be positive: color both passive inputs $r$ and give
the head a foreign singleton. Losing the head would give a
nonunanimous singleton output. Antitonicity gives all subcolorings.
Thus a candidate is specified by a list $Z$ of losses $(j,t)$, each
coloring both other inputs with distinct owners. Call $(r,T)$, with
$\varnothing\ne T\subseteq V$, an allowed fold when
\[
 (j,t)\in Z,\ j\in T\quad\Longrightarrow\quad T\cap t_r\ne\varnothing.
\]
These are precisely the folds $\mathop{\star_r}_{u\in T}u$ omitting
every listed mark. They can be read off the list: a singleton requires
an input with no losses; a pair requires that each head of the pair
sees the other input colored $r$ at every loss; the full set requires
$r$ in every loss.

\begin{proposition}[Two conditions for ternary unanimous guards]
\label{core:prop:ternary-fold-transversals}
For three inputs and any $K\ge3$, $Z$ is the exact list of losses of
a core operation with unanimous singleton profiles if and only if
both conditions hold.
\begin{enumerate}
\item On a row of three distinct singleton inputs, at most one of its
winner omissions is listed.
\item For each positive source query $(i,p)$ there is an allowed
fold $(r,T)$ with $T\cap p_r=\varnothing$.
\end{enumerate}
\end{proposition}
\begin{proof}
The first condition makes the reconstructed operation unanimous. A
loss requires two distinct passive singletons. If the head is ambiguous
then those two passive sources survive, since each sees an ambiguous
input. On a row using at most two singleton colors, a uniquely
supplied winner sees a monochromatic passive coloring, while a winner
with two suppliers has a positive bit at either supplier. On three
distinct singletons the first condition prevents any singleton output.

Apply \cref{core:cor:input-guard-cuts} to a positive $(i,p)$.
An obstruction with $H=V$ makes every allowed fold meet $p_r$:
if $T\cap p_r=\varnothing$, its cut $U=T$ would need a row at a
head of $T$ with $t_r\cap T=\varnothing$. Conversely, in the absence
of such a fold, every nonempty $U\subseteq V\setminus p_r$ has just
such a row. The cut condition thus holds for $H=V$.

No smaller $H$ obstructs. For $H=\{i\}$ a row with internal tails
must agree with $p$ on both passive positions, contradicting positivity.
For $H=\{i,j\}$ the third input must have a fixed color $c$ in $p$
and in every eligible row. Both heads must occur. If a row at $i$
colors $j$ by $s$ and a row at $j$ colors $i$ by $u\ne s$, the row
with $i=u$, $j=s$ and third input $c$ loses two distinct winners.
Hence every eligible row colors the other head by the same $s$.
To pass the owner-$s$ cut on $H$, some vertex of $H$ must already
lie in $p_s$. Only $j$ can do so, making $p$ itself a listed loss
at $i$. This contradiction proves the second condition sufficient;
the cut with $H=V$ proves necessity.
\end{proof}

Six owners can be needed to see incompatibility. Give each of the
three heads one loss, using six distinct colors across the three
passive pairs, and make all other bits positive. There is no allowed
fold: a pair would need a shared color in two rows and the full set
a color in three. Yet on any palette of at most five owners at least
one of the losses disappears. The unused head then supplies an allowed
singleton fold at every owner, and the first condition holds since two
losses never agree at their common passive position. The restriction
to any such palette of size at least three is therefore core.

\paragraph{Three inputs and three owners.}
Suppose that every owner's singleton profile is unanimity. On a
permutation row $\sigma=\sigma_0\sigma_1\sigma_2$, the three inputs
are the singletons $\{\sigma_0\},\{\sigma_1\},\{\sigma_2\}$.
Write $\sigma-r$ for the assertion that the output omits $r$ there.
Each row admits at most one omission, since its output is ambiguous.

\begin{theorem}[Four rules for ternary unanimous profiles]
\label{core:thm:ternary-unanimity}
For three inputs and three owners, the six permutation outputs determine
a core operation with unanimous singleton profiles uniquely. A choice of
at most one omission per row occurs exactly when it satisfies the following
rules under every independent permutation of inputs and owners:
\begin{equation}\label{core:eq:ternary-loss-rules}
\begin{array}{c@{,\quad}c@{,\quad}c@{\qquad\Longrightarrow\qquad}l}
 012-0&021-1&102-0&120-0\\
 012-0&120-0&201-0&021-0\\
 012-0&021-1&120-2&012-2\\
 012-0&021-1&201-0&012-1 .
\end{array}
\end{equation}
The last two lines thus forbid their premise triples. The four rules
are irredundant under relabeling. All four implications hold for
arbitrary core terms, even without the unanimity assumption.
There are $2149$ operations in this class. Every one is a union of
instances of five terms. Write $R(u,v)=u\star_r v$,
$S(u,v)=u\star_s v$ for $r\ne s$, and $A=S(x,S(y,z))$. The terms are
\[
 \begin{gathered}
 R(x,R(y,z)),\qquad R(x,A),\qquad R(x,R(y,S(x,z))),\\
 R(x,S(x,R(z,S(y,z)))),\qquad R(x,R(y,A)),
 \end{gathered}
\]
where $x,y,z$ may be permuted and $r,s$ chosen anew in each summand.
\end{theorem}
\begin{proof}
All monochromatic colorings and their subcolorings are already positive.
The only remaining colorings of the two passive inputs assign them distinct colors. Each is tested by the
unique permutation row putting the missing color at $i$. This proves
uniqueness. Conversely, a core word that passes the nine monochromatic
tests (one per $i,r$) has unanimous profiles. A winning proper
$r$-coalition would still win after every input except a single
$i$ outside the coalition was set to $\{r\}$, violating that test.

For necessity in \eqref{core:eq:ternary-loss-rules}, apply
\cref{core:thm:injective-blockers} to the three omissions on a line
and query the winner on its right. An omission $\sigma-r$ marks the
position of $r$ and colors the other two positions as in $\sigma$.
In the four lines the queried positions are respectively $2,0,2,1$,
each among the marked positions on its left. For each line and each
owner the reachability set is all three positions. For example, in
the first line the owner $0$ starts at positions $0,1$ (the two
omissions of $0$), and the row $021-1$ gives the edge $0\to2$.
Thus none of the three omissions can be removed at any owner, ruling
out the queried winner.

For irredundancy take exactly the three omissions on the left of any
one line. Its conclusion fails. The four premise triples belong to
distinct orbits: the second and third have one and three distinct
missing winners, respectively. In the first line the two rows losing
the same winner agree in one input position; in the fourth they agree
in none. Thus the other three rules hold under every relabeling.

For sufficiency we give the finite enumeration explicitly. Let $P$ be
the $18$ pairs $(\sigma,r)$ of a permutation row and a winner, and
let $Q$ consist of the nine monochromatic tests just used: at $(i,r)$
test the winner $r+1\pmod3$ with that singleton at $i$ and $r$ at
the other two inputs. For every permuted instance $w$ of the five
terms let $B_w\subseteq P\sqcup Q$ be the tests it passes. If $L$ is
a selection of omissions in $P$, set
\[
 B(L)=\bigcup_{w:\,B_w\cap L=\varnothing} B_w.
\]
The table gives the count by $|L|$. The second row tests coverage by
the indicated union, so it also supplies the realizing word.
\[
\begin{array}{c|rrrrrrr}
 |L|&0&1&2&3&4&5&6\\\hline
 \text{satisfies all four rules}&1&18&135&444&729&612&210\\
 \text{satisfies the rules and }B(L)=(P\setminus L)\sqcup Q
   &1&18&135&444&729&612&210
\end{array}
\]
Here is pseudocode for reproducing both rows. The gate evaluation uses
$\{r\}$ if either operand is $\{r\}$, and their union otherwise.
\begin{quote}\small
\begin{tabbing}
\hspace{1em}\=\hspace{1em}\=\hspace{1em}\=\kill
$\mathcal B\gets\{B_w:\ w\text{ one of the five terms, all }6
   \text{ input orders and }6\text{ owner pairs}\}$\\
$\mathcal R\gets\{(C,d): C\Rightarrow d\text{ a line of
   \eqref{core:eq:ternary-loss-rules} in any input and owner order}\}$\\
for $(u_\sigma)_\sigma\in\{\circ,0,1,2\}^{6}$:\\
\> $L\gets\{(\sigma,u_\sigma):u_\sigma\ne\circ\}$\\
\> if every $(C,d)\in\mathcal R$ satisfies $C\not\subseteq L$ or $d\in L$:\\
\>\> increment row one at $|L|$\\
\>\> if $\bigcup\{B\in\mathcal B:B\cap L=\varnothing\}
             =(P\setminus L)\sqcup Q$, increment row two at $|L|$
\end{tabbing}
\end{quote}
There are $93$ distinct sets in $\mathcal B$. Every passing union has
all nine monochromatic bits and exactly the specified permutation
outputs, so the first paragraph gives the claimed profiles and operation.
\end{proof}
